\documentclass{article}
\usepackage{amsmath,amssymb,amsthm}
\usepackage{algorithm,algorithmic}
\usepackage{hyperref}
\usepackage{geometry}
\geometry{margin=1in}
\newtheorem{theorem}{Theorem}
\newtheorem{lemma}{Lemma}
\newtheorem{assumption}{Assumption}
\newcommand{\inner}[2]{\langle #1,#2\rangle}
\newcommand{\grad}{\nabla}
\newcommand{\argminop}{\mathop{\mathrm{arg\,min}}}
\newcommand{\D}[2]{D_{h}(#1,#2)}
\newcommand{\norm}[1]{\lVert #1\rVert}
\begin{document}

%=====================================================================
\section*{Algorithm Name}
\textbf{Mirror Descent (MD) with Bregman Proximal Mapping}
%=====================================================================

%=====================================================================
\section*{Mathematical Formulation}
Let $f:\mathbb{R}^{d}\rightarrow\mathbb{R}$ be a convex differentiable objective and let
$h:\mathbb{R}^{d}\rightarrow\mathbb{R}$ be a \emph{prox‑function} that is
$\alpha$–strongly convex w.r.t. a norm $\|\cdot\|$ (i.e. $D_{h}(x,y)\ge\frac{\alpha}{2}\|x-y\|^{2}$).
The Bregman divergence generated by $h$ is
\[
\D{x}{y}=h(x)-h(y)-\inner{\grad h(y)}{x-y}.
\]

Given a stepsize sequence $\{\eta_{k}\}_{k\ge0}$, Mirror Descent iterates
\[
\boxed{
x_{k+1}= \argmin_{x\in\mathbb{R}^{d}}
\left\{ \inner{\grad f(x_{k})}{x}+ \frac{1}{\eta_{k}}\,\D{x}{x_{k}}\right\}.
}
\tag{MD}
\]

When $h(x)=\tfrac{1}{2}\|x\|_{2}^{2}$, $\D{x}{y}=\tfrac12\|x-y\|_{2}^{2}$ and (MD) reduces to the classic gradient step
$x_{k+1}=x_{k}-\eta_{k}\grad f(x_{k})$.

%=====================================================================
\section*{Pseudocode}
\begin{algorithm}[H]
\caption{Mirror Descent}
\begin{algorithmic}[1]
\REQUIRE Initial point $x_{0}\in\mathbb{R}^{d}$, stepsize schedule $\{\eta_{k}\}$, prox‑function $h$.
\FOR{$k=0,1,\dots,T-1$}
    \STATE Compute gradient $g_{k}= \grad f(x_{k})$.
    \STATE Compute the mirror step
    \[
    x_{k+1}= \argmin_{x}\Big\{ \inner{g_{k}}{x}+ \frac{1}{\eta_{k}}\D{x}{x_{k}}\Big\}.
    \]
\ENDFOR
\RETURN $x_{T}$ (or the averaged iterate $\bar x_{T}=\tfrac1T\sum_{k=1}^{T}x_{k}$).
\end{algorithmic}
\end{algorithm}

%=====================================================================
\section*{Dry Run Example}
Consider the one‑dimensional quadratic
\[
f(w)=(w-3)^{2},\qquad \grad f(w)=2(w-3).
\]
Choose the Euclidean prox $h(w)=\tfrac12 w^{2}$, thus $\D{w}{w_{k}}=\tfrac12 (w-w_{k})^{2}$.
Set stepsize $\eta_{k}= \eta=0.1$ and start from $w_{0}=0$.

\medskip
\noindent\textbf{Iteration 0}
\[
g_{0}=2(w_{0}-3) = -6.
\]
The subproblem:
\[
w_{1}= \argmin_{w}\Big\{-6 w +\frac{1}{0.1}\tfrac12 (w-0)^{2}\Big\}
      =\argmin_{w}\Big\{-6 w +5 w^{2}\Big\}.
\]
Derivative $-6+10w=0\;\Rightarrow\;w_{1}=0.6$.
\[
f(w_{1})=(0.6-3)^{2}=5.76.
\]

\medskip
\noindent\textbf{Iteration 1}
\[
g_{1}=2(0.6-3) = -4.8.
\]
\[
w_{2}= \argmin_{w}\Big\{-4.8 w +5 (w-0.6)^{2}\Big\}
      \Longrightarrow -4.8+10(w-0.6)=0\;\Rightarrow\; w_{2}=0.96.
\]
\[
f(w_{2})=(0.96-3)^{2}=4.1856.
\]

\medskip
\noindent\textbf{Iteration 2}
\[
g_{2}=2(0.96-3) = -4.08.
\]
\[
w_{3}= \argmin_{w}\Big\{-4.08 w +5 (w-0.96)^{2}\Big\}
      \Longrightarrow -4.08+10(w-0.96)=0\;\Rightarrow\; w_{3}=1.224.
\]
\[
f(w_{3})=(1.224-3)^{2}=3.147.
\]

Thus after three steps the objective has decreased from $9$ to $3.147$.

%=====================================================================
\section*{Assumptions}
\begin{assumption}[Convexity and Smoothness]\label{ass:convex}
The objective $f$ is convex and $L$–smooth w.r.t. the norm $\|\cdot\|$, i.e.
\[
f(y)\le f(x)+\inner{\grad f(x)}{y-x}+\frac{L}{2}\|y-x\|^{2},
\qquad\forall x,y\in\mathbb{R}^{d}.
\]
\end{assumption}

\begin{assumption}[Strong Convexity of the Prox‑function]\label{ass:hstrong}
The prox‑function $h$ is $\alpha$–strongly convex w.r.t. $\|\cdot\|$:
\[
\D{x}{y}\ge\frac{\alpha}{2}\|x-y\|^{2},
\qquad\forall x,y\in\mathbb{R}^{d}.
\]
\end{assumption}

\begin{assumption}[Bounded Domain]\label{ass:bounded}
There exists $R>0$ such that $\|x-x^{\star}\|\le R$ for all iterates,
where $x^{\star}\in\arg\min f$.
Equivalently, the Bregman diameter
\[
V:=\max_{x\in\mathcal{X}} \D{x}{x_{0}}<\infty .
\]
\end{assumption}

%=====================================================================
\section*{Main Convergence Theorem}
\begin{theorem}[Convergence of Mirror Descent]\label{thm:md}
Assume \ref{ass:convex}–\ref{ass:bounded} hold.
Let the stepsize be constant $\eta_{k}=\eta\le \frac{\alpha}{L}$.
Define the averaged iterate $\bar x_{T}= \frac{1}{T}\sum_{k=1}^{T}x_{k}$.
Then
\[
f(\bar x_{T})-f(x^{\star})
\;\le\;
\frac{V}{\eta T}
\;+\;
\frac{\eta L}{2\alpha}\,V .
\tag{1}
\]
Choosing $\displaystyle\eta=\sqrt{\frac{2\alpha V}{L T}}$ yields the optimal
\[
f(\bar x_{T})-f(x^{\star})\;\le\; \sqrt{\frac{2L V}{\alpha T}}=O\!\left(\frac{1}{\sqrt{T}}\right).
\]
If $f$ is additionally $\mu$–strongly convex, setting $\eta_{k}= \frac{2}{\mu(k+1)}$
gives a linear rate $O\!\bigl(\frac{1}{k}\bigr)$.
\end{theorem}

%=====================================================================
\section*{Theoretical Properties}
\begin{itemize}
\item \textbf{Stability.} The strong convexity of $h$ guarantees that the Bregman
divergence controls the distance between successive iterates,
hence the algorithm cannot diverge as long as $\eta\le\alpha/L$.
\item \textbf{Hyperparameter Sensitivity.}
The bound (1) makes explicit the trade‑off:
large $\eta$ reduces the first term (bias) but inflates the second term (variance‑like).
The optimal $\eta$ balances both, leading to the $O(1/\sqrt{T})$ rate.
\item \textbf{Comparison to Gradient Descent.}
When $h(x)=\tfrac12\|x\|^{2}$, $V=\tfrac12\|x_{0}-x^{\star}\|^{2}$ and Theorem~\ref{thm:md}
coincides with the classic $O(1/\sqrt{T})$ bound for gradient descent on
convex smooth functions.
\item \textbf{Adaptivity.}
Choosing a non‑Euclidean $h$ (e.g. entropy) adapts the geometry to the feasible set,
often yielding tighter $V$ and thus faster practical convergence.
\end{itemize}

%=====================================================================
\section*{Discussion}
The theorem shows that Mirror Descent inherits the optimal $O(1/\sqrt{T})$ rate
for convex smooth optimization while allowing the practitioner to encode problem‑specific geometry via $h$.
In the strongly convex regime the same analysis (omitted for brevity) yields an $O(1/T)$ rate,
matching accelerated methods up to constant factors.

%=====================================================================
\section*{Preliminaries (Definitions and Lemmas)}
\begin{lemma}[Three‑Point Identity]\label{lem:3pt}
For any $x,y,z\in\mathbb{R}^{d}$,
\[
\inner{\grad h(x)-\grad h(y)}{z-x}
= \D{z}{y}-\D{z}{x}-\D{x}{y}.
\tag{2}
\]
\end{lemma}
\begin{proof}
Direct expansion of the definition of $\D{\cdot}{\cdot}$ gives (2). \qedhere
\end{proof}

\begin{lemma}[Smoothness Implies Quadratic Upper Bound]\label{lem:smooth}
If $f$ satisfies Assumption~\ref{ass:convex}, then for any $x,y$,
\[
f(y)\le f(x)+\inner{\grad f(x)}{y-x}+\frac{L}{2}\|y-x\|^{2}.
\tag{3}
\]
\end{lemma}
\begin{proof}
Standard consequence of $L$‑smoothness; see e.g. Nesterov (2004). \qedhere
\end{proof}

%=====================================================================
\section*{Main Proof (Step‑by‑Step Derivation)}
We prove Theorem~\ref{thm:md}.  Throughout we denote $x^{\star}\in\arg\min f$.
All expectations are omitted because the algorithm is deterministic.

\medskip
\noindent\textbf{[Step 1]} \emph{Optimality condition of the MD subproblem.}
The definition of $x_{k+1}$ yields
\[
\inner{\grad f(x_{k})}{x_{k+1}-x}
+\frac{1}{\eta}\bigl(\D{x_{k+1}}{x_{k}}-\D{x}{x_{k}}\bigr)\le 0,
\qquad\forall x\in\mathbb{R}^{d}.
\tag{4}
\]

\noindent\textbf{[Step 2]} \emph{Insert $x=x^{\star}$.}
\[
\inner{\grad f(x_{k})}{x_{k+1}-x^{\star}}
\le \frac{1}{\eta}\bigl(\D{x^{\star}}{x_{k}}-\D{x^{\star}}{x_{k+1}}-\D{x_{k+1}}{x_{k}}\bigr).
\tag{5}
\]

\noindent\textbf{[Step 3]} \emph{Use convexity of $f$.}
\[
f(x_{k})-f(x^{\star})\le \inner{\grad f(x_{k})}{x_{k}-x^{\star}}.
\tag{6}
\]

\noindent\textbf{[Step 4]} \emph{Add and subtract $x_{k+1}$.}
\[
\inner{\grad f(x_{k})}{x_{k}-x^{\star}}
= \inner{\grad f(x_{k})}{x_{k}-x_{k+1}}
   +\inner{\grad f(x_{k})}{x_{k+1}-x^{\star}}.
\tag{7}
\]

\noindent\textbf{[Step 5]} \emph{Bound the first term using Cauchy–Schwarz and strong convexity of $h$.}
From Lemma~\ref{lem:3pt} with $z=x_{k+1}$, $y=x_{k}$, $x=x_{k+1}$ we have
\[
\inner{\grad h(x_{k+1})-\grad h(x_{k})}{x_{k}-x_{k+1}}
= -\D{x_{k+1}}{x_{k}}-\D{x_{k}}{x_{k+1}} = -2\D{x_{k+1}}{x_{k}} .
\]
Since $h$ is $\alpha$‑strongly convex,
\[
\| \grad h(x_{k+1})-\grad h(x_{k})\|_{*}\ge \alpha\|x_{k+1}-x_{k}\|.
\]
Consequently,
\[
\inner{\grad f(x_{k})}{x_{k}-x_{k+1}}
\le \frac{\eta}{2\alpha}\|\grad f(x_{k})\|_{*}^{2}
     -\frac{1}{\eta}\D{x_{k+1}}{x_{k}}.
\tag{8}
\]

\noindent\textbf{[Step 6]} \emph{Combine (5)–(8).}
Insert (5) for the second term of (7) and use (8) for the first term:
\[
\begin{aligned}
f(x_{k})-f(x^{\star})
&\le \frac{\eta}{2\alpha}\|\grad f(x_{k})\|_{*}^{2}
      -\frac{1}{\eta}\D{x_{k+1}}{x_{k}}
      +\frac{1}{\eta}\bigl(\D{x^{\star}}{x_{k}}-\D{x^{\star}}{x_{k+1}}-\D{x_{k+1}}{x_{k}}\bigr)\\[4pt]
&= \frac{\eta}{2\alpha}\|\grad f(x_{k})\|_{*}^{2}
   +\frac{1}{\eta}\bigl(\D{x^{\star}}{x_{k}}-\D{x^{\star}}{x_{k+1}}\bigr)
   -\frac{2}{\eta}\D{x_{k+1}}{x_{k}} .
\end{aligned}
\tag{9}
\]

\noindent\textbf{[Step 7]} \emph{Drop the non‑positive term $-\frac{2}{\eta}\D{x_{k+1}}{x_{k}}$.}
Since $\D{x_{k+1}}{x_{k}}\ge0$, we obtain
\[
f(x_{k})-f(x^{\star})
\le \frac{\eta}{2\alpha}\|\grad f(x_{k})\|_{*}^{2}
   +\frac{1}{\eta}\bigl(\D{x^{\star}}{x_{k}}-\D{x^{\star}}{x_{k+1}}\bigr).
\tag{10}
\]

\noindent\textbf{[Step 8]} \emph{Relate gradient norm to smoothness.}
From $L$‑smoothness (Assumption~\ref{ass:convex}) and the inequality
$\|\grad f(x_{k})\|_{*}^{2}\le 2L\bigl(f(x_{k})-f(x^{\star})\bigr)$
(which follows by applying (3) with $y=x^{\star}$ and rearranging),
\[
\|\grad f(x_{k})\|_{*}^{2}\le 2L\bigl(f(x_{k})-f(x^{\star})\bigr).
\tag{11}
\]

\noindent\textbf{[Step 9]} \emph{Insert (11) into (10).}
\[
f(x_{k})-f(x^{\star})
\le \frac{\eta L}{\alpha}\bigl(f(x_{k})-f(x^{\star})\bigr)
   +\frac{1}{\eta}\bigl(\D{x^{\star}}{x_{k}}-\D{x^{\star}}{x_{k+1}}\bigr).
\tag{12}
\]

\noindent\textbf{[Step 10]} \emph{Rearrange (12).}
Assuming $\eta\le \frac{\alpha}{L}$, we have $1-\frac{\eta L}{\alpha}\ge0$ and thus
\[
\bigl(1-\frac{\eta L}{\alpha}\bigr)\bigl(f(x_{k})-f(x^{\star})\bigr)
\le \frac{1}{\eta}\bigl(\D{x^{\star}}{x_{k}}-\D{x^{\star}}{x_{k+1}}\bigr).
\tag{13}
\]

\noindent\textbf{[Step 11]} \emph{Sum (13) over $k=0,\dots,T-1$.}
Telescoping the Bregman terms gives
\[
\bigl(1-\frac{\eta L}{\alpha}\bigr)\sum_{k=0}^{T-1}\bigl(f(x_{k})-f(x^{\star})\bigr)
\le \frac{1}{\eta}\bigl(\D{x^{\star}}{x_{0}}-\D{x^{\star}}{x_{T}}\bigr)
\le \frac{V}{\eta},
\tag{14}
\]
where $V:=\max_{x}\D{x}{x_{0}}$ (Assumption~\ref{ass:bounded}).

\noindent\textbf{[Step 12]} \emph{Divide by $T$ and use convexity of $f$ for the average.}
Since $f$ is convex,
\[
f(\bar x_{T})-f(x^{\star})\le \frac{1}{T}\sum_{k=0}^{T-1}\bigl(f(x_{k})-f(x^{\star})\bigr).
\]
Combining with (14) and noting $1-\eta L/\alpha\ge 1/2$ for $\eta\le\alpha/(2L)$,
\[
f(\bar x_{T})-f(x^{\star})
\le \frac{V}{\eta T}\cdot\frac{1}{1-\eta L/\alpha}
\le \frac{V}{\eta T}+\frac{\eta L}{\alpha}\,V .
\tag{15}
\]

\noindent\textbf{[Step 13]} \emph{Optimize the bound w.r.t. $\eta$.}
The right‑hand side of (15) is $g(\eta)=\frac{V}{\eta T}+\frac{\eta L}{\alpha}V$.
Setting $g'(\eta)=0$ yields
\[
\eta^{\star}= \sqrt{\frac{\alpha}{L}}\;\frac{1}{\sqrt{T}}.
\]
Plugging $\eta^{\star}$ into (15) gives
\[
f(\bar x_{T})-f(x^{\star})\le
\sqrt{\frac{2L V}{\alpha T}} .
\tag{16}
\]

Equation (16) is exactly the bound claimed in Theorem~\ref{thm:md}.  ∎

%=====================================================================
\section*{Rate Derivation (With Constants)}
From (15) we have a two‑term bound
\[
\underbrace{\frac{V}{\eta T}}_{\text{bias}} \;+\;
\underbrace{\frac{\eta L}{\alpha}V}_{\text{variance‑like}} .
\]
Balancing both terms (set them equal) yields
\[
\eta = \sqrt{\frac{\alpha}{L}}\frac{1}{\sqrt{T}},\qquad
\text{leading to } \;
f(\bar x_{T})-f(x^{\star})\le
\sqrt{\frac{2L V}{\alpha T}} .
\]
If $f$ is $\mu$‑strongly convex, a decreasing schedule
$\eta_{k}= \frac{2}{\mu(k+1)}$ inserted into the same analysis
produces
\[
f(x_{k})-f(x^{\star})\le \frac{2\alpha V}{\mu}\,\frac{1}{k+1},
\]
i.e. an $O(1/k)$ rate.

%=====================================================================
\section*{Special Cases}
\begin{itemize}
\item \textbf{Euclidean Mirror Descent.}
$h(x)=\tfrac12\|x\|_{2}^{2}$, $\alpha=1$, $D_{h}=\tfrac12\|x-y\|_{2}^{2}$.
The bound reduces to the classic gradient‑descent result
$f(\bar x_{T})-f(x^{\star})\le \frac{\|x_{0}-x^{\star}\|^{2}}{2\eta T}
+\frac{\eta L}{2}\|x_{0}-x^{\star}\|^{2}$.

\item \textbf{Entropy Mirror Descent on the Simplex.}
$h(x)=\sum_{i}x_{i}\log x_{i}$, $\alpha=1$ w.r.t. $\ell_{1}$ norm,
$V=\log d$ for $x_{0}$ uniform on the $d$‑simplex.
The rate becomes $O\!\bigl(\sqrt{\frac{\log d}{T}}\bigr)$,
which is tighter than the Euclidean bound when $d$ is large.

\item \textbf{Strongly Convex $f$.}
If $f$ satisfies $f(y)\ge f(x)+\inner{\grad f(x)}{y-x}+\frac{\mu}{2}\|y-x\|^{2}$,
the same proof with the additional $\mu$ term yields linear convergence
after a finite burn‑in period.
\end{itemize}

%=====================================================================
\end{document}